% Standalone paragraph. Paste into your own thesis, or \input it.
% Not wired into any chapter's numbering -- drop it wherever the "the population is wrong" argument
% is needed. Companion to whole_tree_gap.tex and reproducible_features.tex: those argue the field
% measures the wrong place in the tree and cannot agree on how to measure it, this one argues it
% measures the wrong people.
%
% STANDING CONSTRAINT respected here: no exact figure from the Copenhagen General Population Study
% appears in this file, or anywhere under thesis/, until data access.
% CHECK BEFORE SUBMISSION: the "Herlev-Osterbro" alias for the CGPS is used informally in this
% project; confirm the study's own preferred wording before the parenthetical stays in.

% HEADLINE: one of three sibling subsections for the State of the Art chapter, sitting under
% \section{The gap} (state_of_the_art/08_the_gap.tex): where the tree is measured
% (whole_tree_gap), how it is measured (reproducible_features), in whom it is measured
% (healthy_cohort_2). Promote to \section if they are to stand as sections of their own.
\subsection{In whom the tree is measured}
\label{sec:gap-whom}

Coronary geometry is nearly always measured in people who had a clinical reason to be scanned. For the question the field has mostly asked --- where in the tree plaque appears --- that is an
appropriate population: a lesion has to be present to be located.
It is well documented that lesions are not distributed at random: they form where flow slows or briefly reverses, since low and
oscillatory shear on the vessel lining activates the gene program that initiates atherosclerosis
\cite{chatzizisis2007ess}.
Flow traced through human coronary trees after death localized plaque almost exclusively to the
outer wall of a vessel leaving a bifurcation and the inner wall of a curve \cite{asakura1990flow},
and imaging all three arteries in 131 patients found the rupture-prone phenotype clustered
proximally in each \cite{araki2020plaques}.


The clinical indication is also why such data exists in quantity.
CCTA is the recommended first-line test when chronic coronary syndrome is suspected
\cite{knuuti2019esc}, so those scans accumulate as a by-product of ordinary care, whereas imaging
people who report no symptoms is not part of routine practice and happens mainly inside studies
designed for it.

The problem, for questions like what a normal tree looks like, or which of its features carry risk
before disease develops, is that disease alters the geometry it would be predicted from.
Examining 136 hearts at autopsy, Glagov et al.\ found the artery enlarges outward as plaque
accumulates, so the lumen does not measurably narrow until the lesion occupies roughly 40\% of the
cross-section \cite{glagov1987remodeling}.
An artery called normal in a referred cohort may therefore already have remodeled, so a shape
measured in it can be the disease's product rather than the risk that preceded it.
Normal coronary anatomy has been described before --- 300 scans selected for zero calcium and no
stenosis \cite{medranogracia2016atlas}, 281 patients characterized automatically
\cite{nannini2024characterization} --- but selecting on absent findings is not sampling a
population, and both stop at description.
Three constraints follow, each corresponding to a question this thesis asks: the population is
selected for disease, each participant is imaged once, and the description is never tested against
what happens afterwards.
The field acknowledges the first, its recent review naming case selection and sample size among the
limits that ``may undermine the generalisability of the analyses'' and expecting future studies to
focus on more extensive and diverse populations \cite{shen2026geometry}.


With the Copenhagen General Population Study (Herlev--{\O}sterbro) we may have a real opportunity to
do something that has not been done before. To my understanding these participants were imaged as part of a research protocol rather than because anyone suspected
disease in them, so their trees form the distribution of a population rather than of patients.
A substantial number have also been scanned a second time, which lets us watch how a tree changes,
or does not, in people in whom no disease is found.
Where participants have had events in the meantime, it will be interesting to see in the second scan how their arteries have changed. And to see whether the trees that were healthy at the first scan
and later had an event share features that the trees without events do not.
